% GENERATED FILE. Edit canonical lessons, then run npm run generate:books.
% canonical-source-hash: fnv1a64:5c2d7f2aecdb0a13
% canonical-lessons: GE-C272-sendung, GE-C272-rundfunk, GE-C272-presse, GE-C272-journalist, GE-C272-titelseite

\chapter{Im Radio}
\label{ch:ge-im-radio}
\hlchaptermodality{272}

\begin{chapteropening}
\textbf{By the end of this chapter:} \emph{I can talk about where the news comes from: a broadcast, radio, the press, a journalist and the front page.}

Complete the last lesson of chapter 272: I can talk about where the news comes from: a broadcast, radio, the press, a journalist and the front page.
\end{chapteropening}

\section[die Sendung]{die Sendung — the broadcast, the programme; the consignment}

\label{lesson:GE-C272-sendung}



\begin{quote}
Before the new one: say the German for the headline, then the German for the article.
\end{quote}

\subsection*{You'll want to know: die Sendung}
\textbf{die Sendung} — \textquotedblleft{}the broadcast, the programme; the consignment\textquotedblright{}.

\begin{quote}
Die Sendung beginnt um acht Uhr.
\end{quote}

\begin{quote}
\emph{The programme begins at eight.}
\end{quote}

Plural: \emph{die Sendungen}. At the post office, a \emph{Sendung} is a parcel or letter on its way.

\begin{cousinweb}
\emph{Sendung} is built on \emph{senden} \textquotedblleft{}to send\textquotedblright{}, a cousin of English \emph{send} from Proto-Germanic \emph{sandijaną} \textquotedblleft{}to cause to go\textquotedblright{} (EtymDB).
\end{cousinweb}

\begin{grammarlens}[title={What you've built}]
One more word for telling, explaining and describing at B1.
\end{grammarlens}

\subsection*{Your turn}
\begin{itemize}
\raggedright
  \item \emph{Say it:} \emph{die Sendung}
  \item \emph{Say it:} \emph{die Sendung}, once more
  \item \emph{Say it:} a sentence of your own with \emph{die Sendung}
  \item \emph{Recall:} say the German for the headline, then the German for the article, then say \emph{die Sendung} again
\end{itemize}

\begin{tcolorbox}[breakable,colback=teal!4,colframe=teal!35!black,title={Before you move on}]
What does die Sendung mean? (\textquotedblleft{}The broadcast, the programme; the consignment\textquotedblright{}.) Say it once more.
\end{tcolorbox}
\section[der Rundfunk]{der Rundfunk — radio, broadcasting}

\label{lesson:GE-C272-rundfunk}



\begin{quote}
Before the new one: say the German for the article, then the German for the broadcast.
\end{quote}

\subsection*{You'll want to know: der Rundfunk}
\textbf{der Rundfunk} — \textquotedblleft{}radio, broadcasting\textquotedblright{}.

\begin{quote}
Sein Vater arbeitete viele Jahre beim Rundfunk.
\end{quote}

\begin{quote}
\emph{His father worked in radio for many years.}
\end{quote}

No plural; a little formal. In speech people say \emph{das Radio}.

\begin{cousinweb}
\emph{Rundfunk} joins \emph{rund} \textquotedblleft{}round\textquotedblright{}, Middle High German \emph{runt} (EtymDB), which came through French from Latin \emph{rotundus} \textquotedblleft{}round\textquotedblright{} (Lewis \& Short), and \emph{Funk} \textquotedblleft{}radio\textquotedblright{}, a by-form of \emph{Funke} \textquotedblleft{}spark\textquotedblright{}, from Old High German \emph{funcho} (EtymDB).
\end{cousinweb}

\begin{grammarlens}[title={What you've built}]
One more word for telling, explaining and describing at B1.
\end{grammarlens}

\subsection*{Your turn}
\begin{itemize}
\raggedright
  \item \emph{Say it:} \emph{der Rundfunk}
  \item \emph{Say it:} \emph{der Rundfunk}, once more
  \item \emph{Say it:} a sentence of your own with \emph{der Rundfunk}
  \item \emph{Recall:} say the German for the article, then the German for the broadcast, then say \emph{der Rundfunk} again
\end{itemize}

\begin{tcolorbox}[breakable,colback=teal!4,colframe=teal!35!black,title={Before you move on}]
What does der Rundfunk mean? (\textquotedblleft{}Radio, broadcasting\textquotedblright{}.) Say it once more.
\end{tcolorbox}
\section[die Presse]{die Presse — the press; a squeezer}

\label{lesson:GE-C272-presse}



\begin{quote}
Before the new one: say the German for the broadcast, then the German for radio.
\end{quote}

\subsection*{You'll want to know: die Presse}
\textbf{die Presse} — \textquotedblleft{}the press; a squeezer\textquotedblright{}.

\begin{quote}
Die Presse schrieb viel über den Fall.
\end{quote}

\begin{quote}
\emph{The press wrote a lot about the case.}
\end{quote}

Also a kitchen press: \emph{eine Knoblauchpresse}.

\begin{cousinweb}
\emph{Presse} is built on \emph{pressen} \textquotedblleft{}to press\textquotedblright{}, Old High German \emph{pressōn} (EtymDB), which came from Latin \emph{pressāre}, built on \emph{premere} \textquotedblleft{}to press\textquotedblright{} (Lewis \& Short): first the printing press.
\end{cousinweb}

\begin{grammarlens}[title={What you've built}]
One more word for telling, explaining and describing at B1.
\end{grammarlens}

\subsection*{Your turn}
\begin{itemize}
\raggedright
  \item \emph{Say it:} \emph{die Presse}
  \item \emph{Say it:} \emph{die Presse}, once more
  \item \emph{Say it:} a sentence of your own with \emph{die Presse}
  \item \emph{Recall:} say the German for the broadcast, then the German for radio, then say \emph{die Presse} again
\end{itemize}

\begin{tcolorbox}[breakable,colback=teal!4,colframe=teal!35!black,title={Before you move on}]
What does die Presse mean? (\textquotedblleft{}The press; a squeezer\textquotedblright{}.) Say it once more.
\end{tcolorbox}
\section[der Journalist]{der Journalist — the journalist}

\label{lesson:GE-C272-journalist}



\begin{quote}
Before the new one: say the German for radio, then the German for the press.
\end{quote}

\subsection*{You'll want to know: der Journalist}
\textbf{der Journalist} — \textquotedblleft{}the journalist\textquotedblright{}.

\begin{quote}
Ein Journalist stellte viele Fragen.
\end{quote}

\begin{quote}
\emph{A journalist asked many questions.}
\end{quote}

Weak noun: \emph{den Journalisten}. Feminine \emph{die Journalistin}. The \emph{J} sounds as in French.

\begin{cousinweb}
\emph{Journalist} was borrowed from French \emph{journaliste} (EtymDB), built on \emph{journal} \textquotedblleft{}daily paper\textquotedblright{}, from Latin \emph{diurnālis} \textquotedblleft{}daily\textquotedblright{}, from \emph{diurnus} \textquotedblleft{}of the day\textquotedblright{}.
\end{cousinweb}

\begin{grammarlens}[title={What you've built}]
One more word for telling, explaining and describing at B1.
\end{grammarlens}

\subsection*{Your turn}
\begin{itemize}
\raggedright
  \item \emph{Say it:} \emph{der Journalist}
  \item \emph{Say it:} \emph{der Journalist}, once more
  \item \emph{Say it:} a sentence of your own with \emph{der Journalist}
  \item \emph{Recall:} say the German for radio, then the German for the press, then say \emph{der Journalist} again
\end{itemize}

\begin{tcolorbox}[breakable,colback=teal!4,colframe=teal!35!black,title={Before you move on}]
What does der Journalist mean? (\textquotedblleft{}The journalist\textquotedblright{}.) Say it once more.
\end{tcolorbox}
\section[die Titelseite]{die Titelseite — the front page}

\label{lesson:GE-C272-titelseite}



\begin{quote}
Before the new one: say the German for the press, then the German for the journalist.
\end{quote}

\subsection*{You'll want to know: die Titelseite}
\textbf{die Titelseite} — \textquotedblleft{}the front page\textquotedblright{}.

\begin{quote}
Das Foto war auf der Titelseite.
\end{quote}

\begin{quote}
\emph{The photo was on the front page.}
\end{quote}

Plural: \emph{die Titelseiten}.

\begin{cousinweb}
\emph{Titelseite} joins \emph{Titel} \textquotedblleft{}title\textquotedblright{}, from Latin \emph{titulus} \textquotedblleft{}an inscription, a title\textquotedblright{} (Lewis \& Short), and the inherited \emph{Seite} \textquotedblleft{}side\textquotedblright{}, Old High German \emph{sīta} (EtymDB).
\end{cousinweb}

\begin{grammarlens}[title={What you've built}]
One more word for telling, explaining and describing at B1.
\end{grammarlens}

\subsection*{Your turn}
\begin{itemize}
\raggedright
  \item \emph{Say it:} \emph{die Titelseite}
  \item \emph{Say it:} \emph{die Titelseite}, once more
  \item \emph{Say it:} a sentence of your own with \emph{die Titelseite}
  \item \emph{Recall:} say the German for the press, then the German for the journalist, then say \emph{die Titelseite} again
\end{itemize}

\begin{tcolorbox}[breakable,colback=teal!4,colframe=teal!35!black,title={Before you move on}]
What does die Titelseite mean? (\textquotedblleft{}The front page\textquotedblright{}.) Say it once more.
\end{tcolorbox}
